\chapter{Quadratic Fields and Their Integers}\label{ch:quadratic}

\begin{center}\itshape
The golden ratio satisfies a monic integer polynomial, yet has a denominator $2$.\\
Which fractions deserve to be integers when the number system grows?
\end{center}

The three rings of Chapter~\ref{ch:gaussian} were selected because rounding
made their arithmetic work. We now ask a question that does not depend on a
convenient drawing: inside $\Q(\sqrt d)$, which elements satisfy a monic
integer polynomial? Trace and norm will answer it. For positive $d$,
Pell's equation then supplies infinitely many units; their apparently
irregular sizes become equally spaced after taking logarithms.
The next chapter will use this ring of integers to repair factorization
with ideals.

Throughout this chapter and the next, $d\ne0,1$ is a squarefree integer and $K=\Q(\sqrt d)=\{a+b\sqrt d:a,b\in\Q\}$, a subfield of $\C$ (the inverse of $a+b\sqrt d\ne0$ is $\frac{a-b\sqrt d}{a^2-db^2}$, and $a^2-db^2\ne0$ because $\sqrt d\notin\Q$). The field is \emph{imaginary} if $d<0$ and \emph{real} if $d>0$.



\section{The integers of a quadratic field}

The map $\sigma:a+b\sqrt d\mapsto\overline{a+b\sqrt d}=a-b\sqrt d$ is a field automorphism (a bijection preserving the field operations) of $K$ fixing $\Q$: additivity is clear, and multiplicativity is the identity $(a-b\sqrt d)(a'-b'\sqrt d)=(aa'+dbb')-(ab'+a'b)\sqrt d$. For $d<0$ it is complex conjugation. The \emph{trace} and \emph{norm} of $\alpha=a+b\sqrt d$ are
\[
\operatorname{Tr}\alpha=\alpha+\bar\alpha=2a,\qquad N\alpha=\alpha\bar\alpha=a^2-db^2,
\]
so that $\alpha$ is a root of $X^2-(\operatorname{Tr}\alpha)X+N\alpha$. The trace is additive, the norm is multiplicative, and both are invariant under $\sigma$.

\begin{think}
Is the golden ratio $\varphi=\frac{1+\sqrt5}2$ a root of a monic polynomial with integer coefficients? And $\frac{1+\sqrt{-5}}2$? And $\frac12$? Try before reading on, then guess for which $d$ the number $\frac{1+\sqrt d}2$ deserves to be called an integer.
\end{think}

A complex number is an \emph{algebraic integer} if it is a root of a monic polynomial with integer coefficients. We write $\mathcal O_K$ for the set of algebraic integers of $K$; the question is what it looks like.

\begin{proposition}\label{prop:intcrit}
An element $\alpha\in K$ is an algebraic integer if and only if $\operatorname{Tr}\alpha\in\Z$ and $N\alpha\in\Z$. In particular $\mathcal O_K\cap\Q=\Z$, and $\mathcal O_K$ is stable under $\sigma$.
\end{proposition}
\begin{proof}
If the trace and norm are integers, the monic polynomial
$X^2-(\operatorname{Tr}\alpha)X+N\alpha$ proves integrality.
Conversely, let a monic $P\in\Z[X]$ vanish at $\alpha$.
For rational $\alpha$, Exercise~\ref{ex:sqrt} gives $\alpha\in\Z$,
and its trace and norm are integers.

Suppose $\alpha\notin\Q$ and put
$\mu=X^2-(\operatorname{Tr}\alpha)X+N\alpha$.
Division in $\Q[X]$ gives $P=\mu C+R$ with $\deg R\le1$.
The equality $R(\alpha)=0$ forces $R=0$, since a nonzero linear
rational polynomial has only a rational root. Thus $P=\mu C$, and
both $\mu$ and $C$ are monic.

Here is the precise monic-divisor form of Gauss' lemma that we need.
Choose positive rational numbers $c,e$ so that
$A=c\mu$ and $B=eC$ are primitive integer polynomials: clear
denominators and divide the coefficients by their positive gcd.
A product of primitive integer polynomials is primitive.
Indeed, reducing modulo any prime leaves two nonzero polynomials
over a field, whose product is nonzero. Now $AB=ceP$.
Write $ce=u/v$ in lowest positive terms. From $vAB=uP$,
the leading coefficient gives $v\,\mathrm{lc}(AB)=u$;
thus $v\mid u$ and $v=1$. Since $AB=uP$ is primitive,
$u=1$. Finally $c=\mathrm{lc}(A)$ and $e=\mathrm{lc}(B)$
are positive integers because $\mu,C$ are monic.
Their product is $1$, so $c=e=1$, proving $\mu\in\Z[X]$.

This establishes the integer trace and norm. It also proves
$\mathcal O_K\cap\Q=\Z$; conjugation preserves integrality
because it preserves both quantities.
\end{proof}

\begin{proposition}\label{prop:OK}
Let $\omega_d=\sqrt d$ if $d\equiv2,3\pmod4$, and $\omega_d=\frac{1+\sqrt d}2$ if $d\equiv1\pmod4$. Then
\[
\mathcal O_K=\Z[\omega_d]=\{m+n\omega_d:m,n\in\Z\},
\]
every element being written so in exactly one way. In particular $\mathcal O_K$ is a ring.
\end{proposition}
\begin{proof}
Let $\alpha=a+b\sqrt d\in\mathcal O_K$, so $m=2a$ and $n=a^2-db^2$ are integers by Proposition~\ref{prop:intcrit}. Then $d(2b)^2=m^2-4n\in\Z$. Write $2b=\frac rs$ in lowest terms with $s\ge1$: $s^2\mid dr^2$ and $\gcd(r,s)=1$, so $s^2\mid d$, and $s=1$ since $d$ is squarefree. Hence $a=\frac m2$, $b=\frac k2$ with $m,k\in\Z$, and $n\in\Z$ becomes
\[
m^2-dk^2\equiv0\pmod4.
\]
Squares are $0$ or $1$ modulo $4$, according to parity. If $d\equiv2$ or $3\pmod4$ and $k$ is odd, $m^2\equiv d\pmod4$ is impossible; so $k$ is even, then $m^2\equiv0$, so $m$ is even and $\alpha\in\Z[\sqrt d]$. If $d\equiv1\pmod4$ the condition reads $m^2\equiv k^2\pmod 4$, that is $m\equiv k\pmod2$, and then $\alpha=\frac{m+k\sqrt d}2=\frac{m-k}2+k\omega_d\in\Z[\omega_d]$. Conversely every element of $\Z[\omega_d]$ satisfies the congruence (the computation is reversible), hence lies in $\mathcal O_K$. Uniqueness of $m+n\omega_d$ holds because $1,\sqrt d$ are linearly independent over $\Q$. Finally $\omega_d^2=d$ or $\omega_d^2=\omega_d+\frac{d-1}4$, so $\Z[\omega_d]$ is closed under products.
\end{proof}

So $\mathcal O_K$ is $\Z[i]$ for $d=-1$, $\Z[\omega]$ (Eisenstein, since $\frac{1+\sqrt{-3}}2=1+\omega$) for $d=-3$, $\Z[\sqrt{-2}]$ for $d=-2$: the three rings of Chapter~\ref{ch:gaussian}. For $d=5$ it is $\Z[\varphi]$, $\varphi=\frac{1+\sqrt5}2$ the golden ratio; for $d=-5$ it is $\Z[\sqrt{-5}]$. From now on we write $\mathcal O$ for $\mathcal O_K$ and $\omega$ for $\omega_d$. Divisibility, units, associates and irreducibles are defined in $\mathcal O$ as in Chapter~\ref{ch:gaussian}.

\subsection*{A laboratory for the integrality test}
For $K=\Q(\sqrt5)$, the number $\varphi=(1+\sqrt5)/2$ has trace $1$ and norm
$-1$, so it is an algebraic integer, with polynomial $X^2-X-1$.
By contrast, $(1+\sqrt{-5})/2$ has trace $1$ and norm $3/2$, so it is not
integral. A denominator $2$ is not automatically bad; the norm decides.
For a rational number $r$, trace integrality alone is insufficient:
$r=1/2$ has trace $1$ but norm $1/4$.

There is also a useful computational reading of the norm. In the basis
$1,\omega$, multiplication by $\alpha$ is an integer matrix. For example, when
$\omega=\varphi$ and $\alpha=m+n\varphi$, its columns are the coordinates of
$\alpha$ and $\alpha\varphi$, so the matrix is
\[
\begin{pmatrix}m&n\\ n&m+n\end{pmatrix},\qquad
\det=m^2+mn-n^2=N(m+n\varphi).
\]
A unit gives an invertible integer matrix, hence determinant $\pm1$.
The following proposition proves the converse directly and avoids assuming
anything about matrices over $\Z$.

\section{Units}

\begin{proposition}\label{prop:qunits}
$\varepsilon\in\mathcal O$ is a unit if and only if $N\varepsilon=\pm1$. If $d<0$, the units are $\pm1$, except for $d=-1$ (four units) and $d=-3$ (six units).
\end{proposition}
\begin{proof}
If $\varepsilon\eta=1$ then $N\varepsilon\,N\eta=1$ in $\Z$, so $N\varepsilon=\pm1$; conversely $\varepsilon\cdot(\pm\bar\varepsilon)=1$ with $\pm\bar\varepsilon\in\mathcal O$. Let $d<0$ and $\varepsilon=\frac{m+k\sqrt d}2$ as in the previous proof, so $4N\varepsilon=m^2+|d|k^2=4$. If $k\ne0$ then $|d|\le4$, so $|d|\in\{1,2,3\}$ ($d$ is squarefree); $d=-2$ gives no solution with $k\ne0$, and $d=-1$, $d=-3$ are Proposition~\ref{prop:quadunits}. If $k=0$, $\varepsilon=\pm1$.
\end{proof}

\begin{think}
In $\Z[\sqrt2]$, $(1+\sqrt2)(-1+\sqrt2)=1$, so $1+\sqrt2$ is a unit, and so are its powers. Compute $(1+\sqrt2)^2$, $(1+\sqrt2)^3$ and their norms. Is there a unit strictly between $1$ and $1+\sqrt2$? Draw the units $x+y\sqrt2$ as points $(x,y)$: on which curves do they lie?
\end{think}

\begin{theorem}[Dirichlet's unit theorem, real quadratic case]\label{thm:realunits}
Let $d>0$. There is a least unit $\varepsilon>1$ in $\mathcal O$, the \emph{fundamental unit}, and the units are exactly the $\pm\varepsilon^k$, $k\in\Z$.
\end{theorem}
\begin{proof}
By Theorem~\ref{thm:pell}, Pell's equation $x^2-dy^2=1$ has a solution with $y\ne0$, which gives a unit $x+y\sqrt d\ne\pm1$; one of $\pm(x+y\sqrt d)^{\pm1}$ is $>1$. Now let $u=x+y\sqrt d>1$ be any unit, with $2x,2y\in\Z$. Since $|u\bar u|=1$, we have $|\bar u|<1$, so $x=\frac{u+\bar u}2>0$ and $y\sqrt d=\frac{u-\bar u}2>0$: thus $x,y\ge\frac12$. Below any bound $B$ there are therefore only finitely many units $>1$ (finitely many half-integers $x,y\in\left[\frac12,B\right]$), and a least one, $\varepsilon$, exists. If $u>0$ is a unit, let $k\in\Z$ with $\varepsilon^k\le u<\varepsilon^{k+1}$; then $u\varepsilon^{-k}$ is a unit in $[1,\varepsilon)$, so it equals $1$ by minimality. A negative unit is $-1$ times a positive one.
\end{proof}

\begin{figure}[ht]
\centering
\fig{ch11-units}
\caption{The units $x+y\sqrt2$ of $\Z[\sqrt2]$ with $x,y\ge0$ lie on the two hyperbolas $x^2-2y^2=\pm1$, alternately. Multiplication by $\varepsilon=1+\sqrt2$ is the linear map $(x,y)\mapsto(x+2y,x+y)$, which preserves the pair of hyperbolas and jumps from each unit to the next.}
\label{fig:units}
\end{figure}

\begin{remark}
For a general number field $K$ of degree $n$, with $r_1$ real and $2r_2$ non-real embeddings into $\C$, Dirichlet (1846) proved that the unit group is a finite cyclic group times $\Z^{r_1+r_2-1}$. Real quadratic fields ($r_1=2$, $r_2=0$) give rank $1$, imaginary ones ($r_1=0$, $r_2=1$) rank $0$, in agreement with the two results above. The general theorem is proved with Minkowski's geometry of numbers in Volume~12.
\end{remark}

\subsection*{Why logarithms turn units into steps}
For a positive unit $u$, its conjugate satisfies $|\bar u|=u^{-1}$.
Thus the two real sizes compensate exactly. For $\varepsilon=1+\sqrt2$,
the positive units have ordinary sizes
\[
\ldots,\varepsilon^{-2},\varepsilon^{-1},1,\varepsilon,
\varepsilon^2,\ldots,
\]
but logarithmic sizes $k\log\varepsilon$, equally spaced.
The finite-lattice argument in the proof prevents units from accumulating near
$1$; the least positive logarithmic step then generates every other step.
This is why a statement about a complicated-looking ring becomes a statement
about the integers $k\in\Z$.

Do not confuse the fundamental unit with the smallest Pell solution of norm
$+1$: in $\Z[\sqrt2]$ the fundamental unit $1+\sqrt2$ has norm $-1$, whereas
$3+2\sqrt2=(1+\sqrt2)^2$ is the first unit greater than $1$ of norm $+1$.

\section{Two coordinates, two embeddings}

Writing $\alpha=a+b\sqrt d$ is a coordinate description.
The two embeddings into $\C$ evaluate $\sqrt d$ at either root
of $X^2-d$: they send $\alpha$ to $a+b\sqrt d$ and $a-b\sqrt d$.
Their sum is the trace and their product the norm.
For $d>0$ both are real; one may be large while the other is tiny.
For $d<0$ they are complex conjugates and the norm is nonnegative.

The distinction matters for units. A real unit of norm $-1$ is still
invertible, with inverse $-\bar\alpha$; positivity of the norm was special
to the imaginary rings of Chapter~\ref{ch:gaussian}.
For $d=13$, the element $(1+\sqrt{13})/2$ has integer trace $1$ and norm
$-3$, so it is integral but not a unit. The element
$\varepsilon=(3+\sqrt{13})/2$ has norm $-1$ and is a unit.
By contrast $\sqrt{13}/2$ has norm $-13/4$ and is not integral.

In fact this $\varepsilon$ is fundamental. A unit $u=x+y\sqrt{13}>1$
has positive half-integers $x,y$. If $u<\varepsilon<4$, then
$y\ge1$ would imply $u\ge1/2+\sqrt{13}>4$.
Thus $y=1/2$, and integrality requires $x=m/2$ for a positive odd $m$.
The norm equation becomes $m^2-13=\pm4$; the least possible positive $m$
is $3$, giving precisely $\varepsilon$. This also illustrates why
rounding a decimal approximation to a unit cannot replace the norm equation.

\section{Exercises}

\begin{exercise}[$\star$]\label{ex:ringint}
Describe $\mathcal O_K$ for $d=-3,-7,5,13$. Show that $\Z[\sqrt{-3}]$ does not have unique factorisation, using $4=2\cdot2=(1+\sqrt{-3})(1-\sqrt{-3})$, and explain why the problem disappears in $\mathcal O_K$.
\end{exercise}

\begin{exercise}[$\star$]\label{ex:goldenunit}
Show that the fundamental unit of $\Q(\sqrt5)$ is the golden ratio $\varphi$, and that $\varphi^n=F_{n-1}+F_n\varphi$ (Fibonacci numbers, $F_0=0$, $F_1=1$). Deduce the units of the smaller ring $\Z[\sqrt5]$, and all solutions of $x^2-5y^2=\pm1$.
\end{exercise}

\begin{exercise}[$\star\star$]\label{ex:t1-norm-matrix}
Use the basis $1,\omega_d$ to find the matrix of multiplication by
$\alpha=a+b\omega_d$ in both congruence cases for $d$.
Show that its trace and determinant equal $\operatorname{Tr}\alpha$
and $N\alpha$. Explain why changing the integral basis does not change
these two quantities.
\end{exercise}

\begin{exercise}[$\star\star$]\label{ex:t1-unit-sign}
Let $\varepsilon>1$ be the fundamental unit of a real quadratic field.
Prove that the least unit greater than $1$ of norm $+1$ is
$\varepsilon$ if $N\varepsilon=1$, and $\varepsilon^2$ otherwise.
For $d=13$, deduce the least positive solution of $x^2-13y^2=1$,
taking account of the smaller ring $\Z[\sqrt{13}]$.
\end{exercise}

\section{Hints}

\begin{hint}{ex:ringint}
Only $d$ modulo $4$ matters. For $\Z[\sqrt{-3}]$, write $\frac{1+\sqrt{-3}}2$ in terms of $\omega=e^{2i\pi/3}$.
\end{hint}

\begin{hint}{ex:goldenunit}
The proof of Theorem~\ref{thm:realunits} gives a lower bound for every unit $>1$. For the second part, look at $F_n$ modulo $2$.
\end{hint}

\begin{hint}{ex:t1-norm-matrix}
Compute $\alpha\cdot1$ and $\alpha\omega_d$. Use
$\omega_d^2=d$ or $\omega_d^2=\omega_d+(d-1)/4$.
\end{hint}

\begin{hint}{ex:t1-unit-sign}
The norm of $\varepsilon^k$ is $(N\varepsilon)^k$.
For $d=13$, calculate successive powers and reduce their coefficients
in the basis $1,(1+\sqrt{13})/2$ modulo $2$.
\end{hint}

\section{Solutions}

\begin{solution}{ex:ringint}
\idea Apply Proposition~\ref{prop:OK}: the answer depends only on $d\bmod4$. For $\Z[\sqrt{-3}]$, irreducibility is a norm computation, and the point is that $\frac{1+\sqrt{-3}}2$ is missing from $\Z[\sqrt{-3}]$; adding it creates the unit that makes $2$ and $1+\sqrt{-3}$ associates.
\writeup $-3\equiv-7\equiv5\equiv13\equiv1\pmod4$, so $\mathcal O_K=\Z\left[\frac{1+\sqrt d}2\right]$ in all four cases: $\Z[\omega]$ (Eisenstein, as $\frac{1+\sqrt{-3}}2=1+\omega$), $\Z\left[\frac{1+\sqrt{-7}}2\right]$, $\Z[\varphi]$ and $\Z\left[\frac{1+\sqrt{13}}2\right]$. In $\Z[\sqrt{-3}]$ the norm is $a^2+3b^2$, which never equals $2$; the units are $\pm1$. So $2$ (norm $4$) and $1\pm\sqrt{-3}$ (norm $4$) are irreducible, since a proper factor would have norm $2$; and $1+\sqrt{-3}\ne\pm2$. So $4=2\cdot2=(1+\sqrt{-3})(1-\sqrt{-3})$ are two different factorisations. In $\mathcal O_K=\Z[\omega]$, however, $1+\sqrt{-3}=2(1+\omega)=-2\omega^2$ and $1-\sqrt{-3}=-2\omega$: both factorisations are $2\cdot2$ up to units. The moral is to use \emph{all} algebraic integers of $K$; a smaller ring like $\Z[\sqrt{-3}]$ can fail to factor uniquely for a silly reason.
\end{solution}

\begin{solution}{ex:goldenunit}
\idea The proof of Theorem~\ref{thm:realunits} gives a lower bound for every unit $>1$; check that $\varphi$ attains it. The formula for $\varphi^n$ follows from $\varphi^2=\varphi+1$ by induction, and membership in $\Z[\sqrt5]$ is a parity question on $F_n$.
\writeup $N\varphi=\frac{1-5}4=-1$, so $\varphi$ is a unit. In that proof, a unit $u=x+y\sqrt5>1$ has $x,y\ge\frac12$, so $u\ge\frac{1+\sqrt5}2=\varphi$: hence $\varepsilon=\varphi$, and the units of $\Z[\varphi]$ are $\pm\varphi^n$, $n\in\Z$. Next, $\varphi^1=F_0+F_1\varphi$, and if $\varphi^n=F_{n-1}+F_n\varphi$ then $\varphi^{n+1}=F_{n-1}\varphi+F_n(\varphi+1)=F_n+F_{n+1}\varphi$. As $\varphi=\frac12+\frac12\sqrt5$, $\varphi^n=\left(F_{n-1}+\frac{F_n}2\right)+\frac{F_n}2\sqrt5$ lies in $\Z[\sqrt5]$ if and only if $F_n$ is even. Modulo $2$, the Fibonacci sequence is $0,1,1,0,1,1,\dots$ (period $3$), so $F_n$ is even if and only if $3\mid n$. (For $n<0$ use $\varphi^{-1}=-\bar\varphi$, which reduces to $n>0$.) The units of $\Z[\sqrt5]$ are therefore $\pm\varphi^{3k}=\pm(2+\sqrt5)^k$, and the solutions of $x^2-5y^2=\pm1$ are $x+y\sqrt5=\pm(2+\sqrt5)^k$, $k\in\Z$, with norm $(-1)^k$: $2^2-5=-1$, $9^2-5\cdot4^2=1$, $38^2-5\cdot17^2=-1$, \dots
\end{solution}

\begin{solution}{ex:t1-norm-matrix}
\idea A multiplication matrix records images of basis vectors as columns;
its invariant trace and determinant recover the two conjugate embeddings.
\writeup For $\omega_d=\sqrt d$ the matrix is
$\begin{pmatrix}a&db\\ b&a\end{pmatrix}$, with trace $2a$ and determinant
$a^2-db^2$. Otherwise put $c=(d-1)/4$. The matrix is
$\begin{pmatrix}a&bc\\ b&a+b\end{pmatrix}$, with trace $2a+b$ and determinant
$a^2+ab-cb^2$. Direct conjugation gives the same field trace and norm.
Changing basis conjugates the matrix by the basis-change matrix.
The determinant is invariant by multiplicativity; the trace is invariant
because $\operatorname{tr}(UV)=\operatorname{tr}(VU)$, as follows by
summing the diagonal entries. These statements use only Volume~0's
finite-dimensional linear algebra.
\end{solution}

\begin{solution}{ex:t1-unit-sign}
\idea Norm signs alternate exactly when the fundamental norm is negative.
Pell solutions also require integer, rather than half-integer, coordinates.
\writeup Every positive unit is $\varepsilon^k$, and those greater than
$1$ have $k>0$. Their norms are $(N\varepsilon)^k$, giving the first claim.
For $d=13$, put $\omega=(1+\sqrt{13})/2$, so $\varepsilon=1+\omega$ and
$\omega^2=\omega+3$. Modulo $2$, multiplication by $\varepsilon$
sends coefficient pairs $(a,b)$ to $(a+b,a)$.
Starting at $(1,0)$ the pairs cycle
$(1,0),(1,1),(0,1),(1,0)$.
Now $a+b\omega$ has integer $\sqrt{13}$-coordinates exactly when $b$ is even.
Thus $\varepsilon^k\in\Z[\sqrt{13}]$ exactly when $3\mid k$.
For norm $+1$ we also need $k$ even, so the least positive exponent is $6$.
Since $\varepsilon^3=18+5\sqrt{13}$,
$\varepsilon^6=649+180\sqrt{13}$, giving $(649,180)$.
Any positive Pell solution is a unit greater than $1$ in the same field,
so the minimality follows from the complete unit description.
\end{solution}

\section*{Sources and further reading}
\begin{itemize}
\item D.~A.~Marcus, \emph{Number Fields}, 2nd ed., Springer, 2018.
Quadratic integers, embeddings and units.
\item H.~W.~Lenstra Jr., Solving the Pell equation,
\emph{Notices of the American Mathematical Society} 49 (2002), 182--192.
The connection between approximation and real quadratic units.
\item P.~Samuel, \emph{Algebraic Theory of Numbers}, Dover, 2008.
The broader arithmetic behind this chapter.
\end{itemize}
